\documentclass{permusi}
\usepackage[utf8]{inputenc}
\usepackage[T1]{fontenc}
\usepackage{geometry}
\usepackage{graphicx}
\usepackage{float}
\usepackage{hyperref}
\usepackage[english,indonesian]{babel}

\title{Optimasi Pemrosesan Template LaTeX Menggunakan Pendekatan Berbasis AI}
\author{Academic Support Agent}
\date{\today}

\begin{document}
\maketitle

\begin{abstract}

\end{abstract}

\chapter{PENDAHULUAN}


\chapter{TINJAUAN PUSTAKA}


% Original content (commented out):
% \documentclass{permusi}
% -------------------------------------------------------
% Packages that facilitate double-blind peer review
% Comment 1st line and uncomment 2nd to remove double-blindness
%\usepackage[blind]{anonymize} % Blind
\usepackage{anonymize} % not blind
% --------------------------------------------------------

% English title
\title{Title of the article:}
\subtitle{subtitle of the article}
% Portuguese / Spanish title
\titulo{Título do artigo:}
\subtitulo{Subtítulo do artigo}

\firstauthor{\anonymize{Author 1}}
\firstorcid{\anonymize{https://orcid.org/1XXX-XXXX-XXXX-XXXX}}
\firstinstitution{\anonymize{University or Research Institute, Department 1}}
\firstemail{\anonymize{E-mail address 1}}

\secondauthor{\anonymize{Author 2}}
\secondorcid{\anonymize{https://orcid.org/1XXX-XXXX-XXXX-XXXX}}
\secondinstitution{\anonymize{University or Research Institute, Department 2}}
\secondemail{\anonymize{E-mail address 2}}

\thirdauthor{\anonymize{Author 3}}
\thirdorcid{\anonymize{https://orcid.org/1XXX-XXXX-XXXX-XXXX}}
\thirdinstitution{\anonymize{University or Research Institute, Department 3}}
\thirdemail{\anonymize{E-mail address 3}}

% \fourthauthor{\anonymize{Author 4}}
% \fourthorcid{\anonymize{https://orcid.org/1XXX-XXXX-XXXX-XXXX}}
% \fourthinstitution{\anonymize{University or Research Institute, Department 4}}
% \fourthemail{\anonymize{E-mail address 4}}

% \fifthauthor{\anonymize{Author 5}}
% \fifthorcid{\anonymize{https://orcid.org/1XXX-XXXX-XXXX-XXXX}}
% \fifthinstitution{\anonymize{University or Research Institute, Department 5}}
% \fifthemail{\anonymize{E-mail address 5}}

% \sixthauthor{\anonymize{Author 6}}
% \sixthorcid{\anonymize{https://orcid.org/1XXX-XXXX-XXXX-XXXX}}
% \sixthinstitution{\anonymize{University or Research Institute, Department 6}}
% \sixthemail{\anonymize{E-mail address 6}}


% -----------------------------------------------
% Header data
% -----------------------------------------------
\runningauthor{Author 1 and Author 2}
\runningtitle{It must fit in one line}
\runningyear{2020}
\runningedition{Computing Performed Music}
\runningnumber{40}
\runninge{e000000}
\runningdoi{XXXXXXXXXXXXXXX}% To be informed
\submitted{XX XXX XXXX} % Format dd Mon YEAR
%\submitted{17 Jan 2020}
\approved{XX XXX XXXX}% To be informed
%\approved{15 Mar 2020}% To be informed


% --------------------------------------------------
% --------------------------------------------------
\begin{document}

\begin{abstract}
 Abstract of no more than 150 words. The text should be 1,15 spacing, 11-point, Calibri, justificed. Text text
text text text text text text text text text text text text text text text text text text text text text text text text text text
text text text text text text text text text text text text text text text text text text text text text text text text text text
text text text text text text text text text text text text text text text text text text text text text text text text text text
text text text text text text text text text text text text text text text text text text text text text text text text text text
text text text text text text text text text text text text text text text text text text text text text text text text text text
text text.
\end{abstract}
\keywords{Keyword 1; keyword 2; keyword 3; keyword 4; keyword 5.}

\begin{resumo}
 Abstract of no more than 150 words. The text should be 1,15 spacing, 11-point, Calibri, justificed. Text text
text text text text text text text text text text text text text text text text text text text text text text text text text text
text text text text text text text text text text text text text text text text text text text text text text text text text text
text text text text text text text text text text text text text text text text text text text text text text text text text text
text text text text text text text text text text text text text text text text text text text text text text text text text text
text text text text text text text text text text text text text text text text text text text text text text text text text text
text text.
\end{resumo}
\palavraschave{Palavra-chave 1; Palavra-chave 2; Palavra-chave 3; Palavra-chave 4; Palavra-chave 5.}


\maketitle

% --------------------------------------------------
\section{Introduction (Use Calibre 20-point)}

Use Calibri, 12-point for the text with double space. 1.15 spacing should be used in all instances, after
periods, colons, semicolons. Add line space before paragraph without indent. The text should be justified.
(Author Year, page). Generic titles of works should be given in \textit{italic}. Flats, sharps, and naturals can be
indicated by the conventional symbols or spelled out (e.g., C-sharp, D-flat, or B-natural). \textbf{Text} text text text
text text text text text text text text text text text text text text text text text text text text text text text text
text text text text text text text text text text text text text text text text text text text text text text text text
text text text text text text text text text text text1. 

\begin{quote}
For citations, use Calibre, 12-point, 1.15 spacing and indent 3.0. Text text text text text text
text text text text text text text text text text text text text text text text text text text text
text text text text text text text text text text text text text text text text text text text text
text text text text text text text text text text text text text text text text text text text text
text text text text text text text text text text text text text text text text text text text text
text text text text text text text text.
\end{quote}

Table \ref{tab:my_label} presents. \lipsum[1]


\begin{table} [htbp]
\caption{A table caption on the top of the table\cite[p22]{butler2011gender}}
\centering
    \begin{tabular}{|l|c|}
        \hline
        \rowcolor{table-title}
        A & Lorem ipsum \\
        \hline
        A1   & Lorem ipsum\\
        A1   & Lorem ipsum\\
        A1   & Lorem ipsum\\
        A1   & Lorem ipsum\\
        A1   & Lorem ipsum\\
        \hline
    \end{tabular}
    \label{tab:my_label}
\end{table}

\lipsum[2-3]

% -----------------------------------------------------
\section{The same thing}

\lipsum[1-10]\footnote{Foot note: Calibri, 12-point, simple spacing, justified.}

\begin{figure}[htpb!]
    \centering
    \includegraphics[width=8cm]{template/logo_permusi.png}    \caption{Caption of the Figure~\cite{becker1961boys}.}
    \label{fig:other_label}
\end{figure}

Figure \ref{fig:other_label} presents \cite{bourdieu1988homo}. \lipsum[1-4]

\subsection{Another subsection}

\lipsum[6-9] \cite{becker1982art}.

Citation example~\cite[p 1]{buscatto2010professional} and another \citep{bruni2000all} here ~\cite{becker1961boys, bruni2000all}, \cite{bourdieu1990domination} 

\bibliographystyle{permusi}
\bibliography{references}

\end{document}


\end{document}